\chapter{主动学习}
\label{chap:active}
\begin{quote}\small
\textit{本章摘要。}主动学习在有限标注预算下选择最有价值的查询。本章从候选池与标注者的交互规则出发，比较不确定性、委员会分歧、预期模型变化和批量多样性，再讨论成本、噪声与流式查询。核心不是找到最难的样本，而是使下一次标注改善部署风险。
\end{quote}

上一章讨论在新数据到来后如何更新而不遗忘，但这通常默认训练所需标签能够取得。工业现场中，更稀缺的资源可能是能够判断裂纹性质、核对工单因果或识别异常波形的专家时间。面对成千上万条未标注记录，更新算法之前还有一个决策：先看哪一条，为什么值得看，得到答案后能够改变什么？本章把这种决策本身纳入学习过程：标注预算是有限资源，查询规则的目标是降低部署风险，而不是让训练集看起来更大。

主动学习不是把最难样本交给专家，而是在当前信息和预算下选择预期有用的观测。模型拿不准可能源于知识不足，也可能源于输入本身无法判读；两个都很有价值的单点放在同一批次中，也可能提供重复信息。因此，本章依次区分预测不确定性与可减少的不确定性，推导标签带来的信息增益和预期风险变化，再用几何覆盖与梯度表示处理批量冗余，最后把费用、标注错误和抽样偏差纳入查询与标注流程。

获取函数给记录排出的次序，只是在估计先标哪一条可能更有帮助。它的分数再高，也要通过实际标注和重新训练来检验。比较方法时应给它们相同的总预算，并计入专家核对、返工和等待时间，再看哪一种更快减少漏报与误报。如果记录分散在不能交换原始数据的工厂中，还需要安排各处怎样共同训练，这就是下一章联邦学习的问题。

\section{查询规则与标注预算}

% auxiliary-resource-guide
本章末的“辅助学习材料”列出与核心方法对应的讲解和实践入口，可在阅读相关推导后，按所列小节对照学习。

\begin{figure}[htbp]
\centering
\begin{tikzpicture}[node distance=0.7cm and 0.75cm]
 \node[idi input, minimum width=1.8cm] (pool) {未标注池\\$U$};
 \node[idi process, right=of pool, minimum width=2.0cm] (score) {获取函数\\$a(x;L)$};
 \node[idi state, right=of score, minimum width=1.8cm] (query) {查询集\\$Q$};
 \node[idi output, right=of query, minimum width=1.8cm] (label) {专家标签};
 \node[idi process, below=0.8cm of score, minimum width=2.2cm] (train) {更新模型\\加入有效标注对};
 \draw[idi arrow] (pool)--(score); \draw[idi arrow] (score)--(query); \draw[idi arrow] (query)--(label); \draw[idi arrow] (label)--(train);
 \draw[idi feedback] (train) to[bend right=35] (score);
\end{tikzpicture}
\caption{池式主动学习在每轮用获取函数分配有限标注预算，再用新标签更新模型和下一轮查询。}
\label{fig:active-query-loop}
\end{figure}
图\ref{fig:active-query-loop}的反馈路径说明查询与训练相互影响：获取函数选择待标注样本，新标签改变模型，也改变下一轮的查询分布。
设已标注集为$L$，未标注池为$U$，标注成本为$c(x)$，总预算为$B$。每轮训练模型，按获取函数$a(x;L)$选择查询集$Q\subset U$，获取标签后更新$L\leftarrow L\cup\{(x,y):x\in Q\}$、$U\leftarrow U\setminus Q$。查询满足$\sum_{x\in Q}c(x)\leq B_{remaining}$。获取函数（acquisition function）是对未来收益的近似，不是测试集上的真实增益\cite{settles2009active}。

池式主动学习可比较整个$U$；流式方法只能对到达样本决定查询或跳过；查询合成允许模型构造新输入，但工业中未必存在可由专家可靠判读的合成波形。无论采用哪种设定，测试集都不属于候选池。初始化标签过少时，模型置信度尚不可靠，因此随机或分层覆盖的种子集十分重要。

\paragraph{一次查询究竟买到了什么。}
一张巡检图像或一段振动窗口是输入$x\in\R^d$，故障类别$y\in\{1,\ldots,K\}$是标签；未标注不是没有输入，而是没有可供训练的可靠$y$。设备是输入的来源，不是一个类别，也不是一次查询。任务是要解决的诊断问题；本章可以始终只有一个任务，却进行很多轮查询。$L$存样本对，$U$只存输入，$Q$在送标注时也只含输入；因此 $|L|$ 计已返回的有效标签，而不计待处理记录。若初始$|L|=4,|U|=10$，本轮查询两条并都得到有效标签，则下一轮为$|L|=6,|U|=8$；标签尚未返回的记录应标为待处理，避免重复派单。拒标仍可能消耗费用，但不能伪造标签加入$L$。

\paragraph{区分训练、选择和验收。}
训练把$L$变成参数$\theta$；获取函数把每个候选变成标量分数；专家查询才产生真实标签；重新训练后独立测试才衡量收益。获取分数不是标签，也不能直接当作训练目标。预算单位应先固定，例如专家分钟数：剩余$10$分钟、每条$3$分钟时最多查询三条，而不是十条。停止条件可以是预算用完或预先规定的验证收益不足，不能反复看最终测试集来决定何时停止。

\section{不确定性、委员会与贝叶斯查询}
\paragraph{最小置信度、间隔与熵。}
设模型类别概率按大小排列为$p_{(1)}(x)\geq p_{(2)}(x)$。常见获取函数为
\begin{align}
 a_{LC}(x)&=1-p_{(1)}(x),\qquad
 a_{margin}(x)=-(p_{(1)}(x)-p_{(2)}(x)),\\
 a_H(x)&=-\sum_kp_k(x)\log p_k(x).
\end{align}
三者均选择得分较大的样本。最小置信度只看最大类别，间隔看前两类竞争，熵考虑全部概率。多分类下它们并不等价。Lewis 与 Gale 的不确定性采样工作展示了基于当前分类器选择标注对象的基本思想\cite{lewis1994uncertainty}。

\paragraph{何时这些分数等价，何时不等价。}
对二分类，记$p=P(Y=1\mid x)$，则最小置信度为$\min(p,1-p)$，负间隔为$-|2p-1|$，熵为$h(p)=-p\log p-(1-p)\log(1-p)$。在$0<p<1/2$时，$h'(p)=\log((1-p)/p)>0$，三者都随$p$靠近$1/2$而增加，故产生相同排序。多分类没有这个结论。例如两个候选概率为$(0.50,0.49,0.01)$与$(0.50,0.25,0.25)$，最小置信度相同，间隔优先选前者，熵却优先选后者；前者突出两类决策竞争，后者突出整体分散程度。这里及下文取自然对数，约定$0\log0=0$。

边界附近可能有可学习样本，也可能是永久模糊的记录、传感器损坏或标签定义冲突。反复查询后一类样本会消耗预算却不减少模型不确定性。应按设备与工况审计获取分数，并检查模型校准，而不是默认低置信度必有高信息价值。

\paragraph{分数的数值和含义。}
二分类候选甲预测为$(0.9,0.1)$，乙为$(0.5,0.5)$，则最小置信度分别为$0.1,0.5$，负间隔为$-0.8,0$，熵约为$0.325,0.693$，三种规则都优先乙。负间隔越接近零越不确定，不能因为它有负号而改成选最小值。熵$H(p)$就是按概率加权的意外程度$-\log p_k$；这里的自然对数使单位为 nat，换成以$2$为底只改变单位和统一缩放，不改变单点排序。预测概率来自当前模型而非真实频率；自信但错误的模型可能永远不主动查询甲。

\paragraph{委员会查询（Query-by-Committee）。}
委员会方法训练多个与当前数据相容的预测器，选择其分歧最大的样本\cite{seung1992qbc}。若有$M$个成员，类别投票比例为$v_k(x)$，投票熵为$-\sum_kv_k\log v_k$。也可用成员概率分布到平均预测的 KL 散度。Bootstrap、不同初始化或后验样本可形成委员会，但高度相关的成员会低估分歧，毫无约束的弱模型则可能制造无效分歧。

\paragraph{Bayesian Active Learning by Disagreement (BALD) 与联合批量查询。}
BALD 选择标签与模型参数间条件互信息最大的输入\cite{houlsby2011bald}：
\begin{equation}
 a_{BALD}(x)=H\!\left[\E_{p(\theta\mid L)}p(y\mid x,\theta)\right]
 -\E_{p(\theta\mid L)}H[p(y\mid x,\theta)].
\end{equation}
第一项是总体预测熵，第二项扣除各模型自身仍有的不确定性。用$M$次后验近似采样计算类别概率、先平均再取熵，并减去逐次熵的平均。它试图区分认知不确定性与数据歧义，但效果取决于后验近似。Monte Carlo (MC) dropout 是一种近似手段，不是精确贝叶斯采样\cite{gal2017deepactive}。

\paragraph{先把两种平均分开计算。}
两个模型分别输出$(0.9,0.1)$和$(0.1,0.9)$。先平均概率得到$(0.5,0.5)$，熵为$0.693$；各自先算熵再平均得到$0.325$，相减为$0.368$。若两个模型都输出$(0.5,0.5)$，相减为零。前者是模型意见相左，标签可能帮助选择模型；后者是每个模型自身都模糊。贝叶斯后验$p(\theta\mid L)$表示看到$L$以后对参数的分布，$\E_{p(\theta\mid L)}$指对不同可能参数求平均，而非对候选样本求平均。认知不确定性指可望随证据减少的参数分歧，数据歧义指给定解释后仍可能出现多个标签；这种分解始终依赖所选模型，不能当作对现实噪声来源的直接测量。

\paragraph{从后验信息增益推出 BALD。}
固定已标注集$L$和候选$x$，把参数$\Theta$看成后验随机变量、标签$Y$看成尚未观测的离散变量。理想查询的信息价值是观察标签后后验改变多少：
\begin{align}
 I(Y;\Theta\mid x,L)
 &=\E_{Y\mid x,L}\KL(p(\theta\mid L,x,Y)\|p(\theta\mid L))\label{eq:active-posterior-information}\\
 &=\E_{\Theta\mid L}\sum_yp(y\mid x,\Theta)
 \log\frac{p(y\mid x,\Theta)}{p(y\mid x,L)}\label{eq:active-predictive-information}\\
 &=H[p(y\mid x,L)]-\E_{\Theta\mid L}H[p(y\mid x,\Theta)].
\end{align}
从式\eqref{eq:active-posterior-information}到式\eqref{eq:active-predictive-information}使用贝叶斯公式$p(\theta\mid L,x,y)/p(\theta\mid L)=p(y\mid x,\theta)/p(y\mid x,L)$，这里采用条件预测模型并将候选$x$作为固定设计。互信息非负且至多为$\log K$，其中$K$为类别数，因为条件熵非负且$K$类熵至多为$\log K$。

取$M$个后验近似样本，令$p^{(m)}\in\R^K$为其预测、$\bar p=M^{-1}\sum_mp^{(m)}$，则估计为
\begin{equation}
 \widehat I=H(\bar p)-\frac1M\sum_mH(p^{(m)})
 =\frac1M\sum_m\KL(p^{(m)}\|\bar p).
\end{equation}
最后一个等式也说明概率委员会分歧与 BALD 的联系，但普通随机初始化委员会未必代表后验。二分类中，若每个模型都输出$(1/2,1/2)$，预测熵为$\log2$而互信息为零：再看标签未必能解决输入本身歧义。若一半模型输出$(1,0)$、另一半输出$(0,1)$，平均预测同样为$(1/2,1/2)$，但每个模型熵为零，互信息达到$\log2$：标签可以区分两组假设。这是两种不确定性的理想示例，不保证近似后验能在真实数据中正确区分它们。

逐点取 BALD 前若干名可能选出几乎相同的相邻窗口。BatchBALD 对一批标签与参数的联合互信息进行近似优化，考虑查询间冗余\cite{kirsch2019batchbald}。批量规模增加使联合计算更困难，因此应报告候选池大小、后验采样次数和获取时间。

对查询批次$Q=\{x_1,\ldots,x_b\}$，若给定参数后标签条件独立，联合信息增益可写为
\begin{equation}
 I(Y_Q;\Theta\mid L)=H(Y_Q\mid L)-\sum_{j=1}^b\E_{\Theta\mid L}H(Y_j\mid x_j,\Theta).
\end{equation}
不能把第一项换成单点熵之和，因为对参数积分后各标签往往相关。两点时，条件独立假设给出$I(Y_1,Y_2;\Theta)=I(Y_1;\Theta)+I(Y_2;\Theta)-I(Y_1;Y_2)$，省略的条件均为固定输入与$L$。这正是冗余扣除：两条记录若几乎回答同一个参数问题，单点相加会重复计算收益。联合标签空间有$K^b$种配置，直接枚举随批量指数增长，所以近似采样和贪心批量构造是计算折中，而不是改变查询目标。

\section{预期风险、梯度变化与多样性}
\paragraph{预期误差减少。}
理想查询希望标注后风险下降。由于未知$y$，可按当前预测枚举候选标签并模拟更新\cite{roy2001error}：
\begin{equation}
 a_{EER}(x)=\widehat R(\theta)
 -\sum_yp_\theta(y\mid x)\widehat R(\theta_{L\cup\{(x,y)\}}).
\end{equation}
风险代理只能在训练侧保留集或规定的无标签代理上计算，不能使用最终测试标签。每个候选、每个标签都重训代价很高，可用少量梯度步近似，但近似可能误判长期收益。若当前模型给真实类别极低概率，预期收益也可能被低估。

\paragraph{怎样模拟还没拿到的标签。}
假设某候选当前预测两类概率为$0.75,0.25$，当前代理风险为$0.20$。分别暂设标签为第一类和第二类，复制模型后独立模拟更新，得到代理风险$0.16,0.24$，则期望更新后风险为$0.75\times0.16+0.25\times0.24=0.18$，获取分数为$0.02$。这两个假想标签都不能写入正式训练集；真正查询后只用专家返回的标签更新。风险代理是训练侧可计算的近似风险，不是预知未来测试误差的工具。若没有代表性验证标签，必须说明无标签代理如何定义，不能把高置信度直接等同于低真实风险。

\paragraph{一次梯度更新能否预测风险收益。}
令$\theta\in\R^p$，候选标签为$y$时的梯度为$g_y(x)=\nabla_\theta\ell(x,y)$，模拟一步$\theta'_y=\theta-\eta g_y$。若风险代理$\widehat R$二阶可微，在当前参数处梯度为$r$、Hessian 为$H$，则
\begin{equation}
 \widehat R(\theta)-\widehat R(\theta'_y)
 \approx\eta r^\top g_y-\frac{\eta^2}{2}g_y^\top Hg_y.
\end{equation}
所以有效更新取决于方向对齐和曲率，并非只取决于$\|g_y\|$。若代理风险在当前模型处已近驻点，一阶项很小，大梯度甚至可能主要增加二阶风险。

还有一个容易忽略的抵消现象。对交叉熵$\ell=-\log p_\theta(y\mid x)$，若用同一个当前模型对未知标签求平均，则
\begin{equation}
 \sum_yp_\theta(y\mid x)g_y(x)
 =-\sum_y\nabla_\theta p_\theta(y\mid x)
 =-\nabla_\theta1=0.
\end{equation}
这里先对每个假想标签的损失求梯度，再用当前概率加权；不是对$\sum_yp_\theta(y\mid x)\ell(x,y)$整体求导，后者还会对概率权重求导。因此简单的一步、固定代理风险、一阶近似组合会给所有候选零平均收益；它不是完整 Expected Error Reduction (EER) 的替代。完整重训会改变最优解、后验或风险估计，高阶项也不会一般消失。这个结果解释了为什么预期梯度长度使用范数的平均，而不是平均梯度的范数，二者由三角不等式并不相同。

预期梯度长度以$\sum_yp_\theta(y\mid x)\|\nabla_\theta\ell(x,y)\|$估计模型变化。较大梯度说明会产生较大更新，却不保证更新方向有益，异常值可能占据排名前列。

\paragraph{Core-set 与 BADGE。}
BADGE 的英文全称为 Batch Active Learning by Diverse Gradient Embeddings，即利用多样化梯度嵌入进行批量主动学习。
Core-set 在表示空间选择覆盖候选池的样本\cite{sener2018coreset}：
\begin{equation}
 \min_{Q\subseteq U:\,|Q|=b}\max_{x\in U}\min_{z\in L_X\cup Q}
 \|f(x)-f(z)\|_2.
\end{equation}
这里$L_X=\{x:(x,y)\in L\}$为已标注集的输入集合，$b$为本轮查询数，$f$为当前特征编码器。这样距离比较的两端都是输入，而不是将带标签样本对与输入混合。
常用贪心近似每次选离已选中心最远的点，并更新所有候选的最近距离。它鼓励覆盖，但若表示由背景或设备身份主导，就会覆盖无关差异；极端离群点也可能被优先选择。

\paragraph{覆盖半径为什么有意义，又为什么不够。}
记当前已选中心集合为$S$，所有候选到中心的最近距离为$d_S(x)$。贪心最远点步骤选择$x^*=\argmax_xd_S(x)$，随后仅需更新$d_{S\cup\{x^*\}}(x)=\min\{d_S(x),\|f(x)-f(x^*)\|\}$，而不必每轮重新比较所有中心。在固定度量、等成本、中心从候选中选取的离散覆盖问题中，最远点策略有经典的两倍半径近似解释：若最终半径超过最优半径的两倍，依次选出的点和最终最远点必须彼此相距超过该值；但最优的$b$个新增半径球无法容纳$b+1$个这样的点，因为同一球内任意两点的距离至多为两倍半径。已有固定中心时，这些点也不可能由已有中心在最优半径内覆盖，故相同计数逻辑仍成立。

这个结论只保证几何覆盖。若条件期望损失$\bar\ell(z)=\E[\ell\mid Z=z]$在表示空间是$L_\ell$-Lipschitz，其中$L_\ell\geq0$是常数而非已标注集$L$，则$|\bar\ell(z)-\bar\ell(z_c)|\leq L_\ell\|z-z_c\|$，覆盖半径$r$才能被解释为至多$L_\ell r$的局部损失差异控制；但故障边界可能不平滑，表示可能遗漏关键变量，而且一个中心标签不等于精确的条件期望损失。因此几何近似保证不能直接改写为标注后分类风险保证。

BADGE 将不确定性与多样性放到最后一层的梯度表示中\cite{ash2020badge}。先取预测概率最大的类别为伪标签，再计算各类别的梯度块：
\begin{equation}
 \hat y=\arg\max_kp_k(x),\qquad
 g_k(x)=(p_k(x)-\mathbf1[k=\hat y])h(x),\label{eq:active-badge-embedding}
\end{equation}
其中$h(x)$是分类头输入。拼接各类梯度后，用类似$k$-means++的距离采样选择一批点。梯度长度反映当前预测的不确定性，梯度方向帮助减少冗余。该方法不需要候选真实标签，但依赖伪标签与当前表示质量。

\paragraph{从覆盖与梯度各算一个小例子。}
在一维表示上，已标注中心为$0$，候选为$1,2,8,9$，预算为两条。最远点法先选$9$，最近中心距离更新为$1,2,1,0$，再选$2$，而不是继续选靠近$9$的$8$；这就是批量多样性的作用。若$9$只是坏传感器的读数，则覆盖仍可能浪费一次查询，故离群审核不能省略。BADGE 中若$h=(1,2)^{\mathsf T}$、$p=(0.6,0.4)$，伪标签为第一类，两个梯度块为$(-0.4,-0.8)^{\mathsf T}$和$(0.4,0.8)^{\mathsf T}$，拼成四维向量。伪标签仅用于算选择特征，不表示已经获得专家标签。

\paragraph{梯度嵌入为何同时包含置信度与方向。}
令线性分类头$W\in\R^{K\times d_h}$、特征$h(x)\in\R^{d_h}$，其中$d_h$是编码后的特征数，不必等于输入维数$d$；logits 为$Wh\in\R^K$，概率为$p\in\R^K$。对标签$y$，交叉熵梯度为
\begin{equation}
 \nabla_W\ell=(p-e_y)h^\top\in\R^{K\times d_h},
\end{equation}
其中$e_y\in\R^K$为第$y$个标准基向量。逐行展平这张矩阵得到$Kd_h$维嵌入，便得到式\eqref{eq:active-badge-embedding}中的$g_k$。其范数满足
\begin{equation}
 \|g(x)\|_2^2=\|p-e_{\hat y}\|_2^2\|h(x)\|_2^2.
\end{equation}
若预测接近某个确定类别，第一个因子接近零；若多个类别竞争，则残差通常更大。但第二个因子也会影响长度，所以不能把它当作纯粹的概率不确定性。两个样本的梯度内积为$(p-e_{\hat y})^\top(p'-e_{\hat y'})\,h^\top h'$：类别竞争结构与特征相似性共同决定方向。

一种可复现的初始化约定是先选梯度范数最大的候选作为首中心，并列时按固定种子的随机规则选择。随后类似$k$-means++，在已有中心集合$S$下，以$D(x)^2/\sum_{u\in U\setminus S}D(u)^2$从未选候选中抽取新中心，其中$D(x)=\min_{z\in S}\|g(x)-g(z)\|$。每选一条即更新最近距离，直到达到本轮预算。这既避免单纯选取最大梯度的一簇近重复样本，又保留随机探索。若所有剩余距离为零，可在未选候选中均匀无放回抽样，不能除以零或重复选择同一条。它是在梯度空间中寻找代表性，并不保证伪标签正确或最终风险最小。

\section{成本、噪声与工业流式变体}
等成本预算下按获取分数排序即可形成简单策略；成本差异明显时，对非负收益分数且$c(x)>0$的情形，可用$a(x)/c(x)$作启发式，或直接求带预算约束的批量选择。前面的负间隔不能未经处理就按此解释为收益：例如固定加一个常数虽不改变单点排序，却会改变除以不同成本后的排序，所以需先明确收益代理的零点与尺度。比值排序不是一般情况下的最优解。长工单、复杂图像与短波形的标注时间不同，应使用专家分钟数或费用，而不只用样本数比较。

若先把收益近似为可加的$a_i\geq0$，批量选择就是$\max_{q_i\in\{0,1\}}\sum_iq_ia_i$，约束$\sum_iq_ic_i\leq B$。比值排序解的是可分割背包的思想，而样本标签不可分割。例如预算为$10$，一个候选成本$6$、收益$12$，另外两个各成本$5$、收益$9$；按比值会先选第一个而得到$12$，选后两个却得到$18$。数值只是算法反例，不是实验结果。实际批量收益还因冗余不满足可加性，故应把多样性和成本联合考虑。

多标注者场景还需选择由谁标注，并允许“不确定”或“无法判读”。对高风险少数故障可配置双人复核；对重复分歧应先修改标签定义。密度加权可以降低孤立噪声的优先级，但也可能排斥真正罕见的新故障，因此应保留随机探索和新颖性审核预算。

流式方法可在$a(x_t)>\tau_t$时查询，并根据剩余预算调整阈值。漂移后旧阈值可能失效，不能仅依赖过去的分数分布。持续学习决定如何更新模型，主动学习决定获取哪些标签，二者可联合但应分别消融。查询高不确定样本产生选择偏差，性能监测仍需要独立的代表性抽样，不能直接用已查询数据的准确率估计上线风险。

\paragraph{查询偏差与代表性监测。}
设一个固定候选总体有$N$个样本，$S_i\in\{0,1\}$表示是否被查询，已知边际入选概率为$\pi_i>0$。对预先固定的评估器，若损失为$\ell_i$，则逆概率估计
\begin{equation}
 \widehat R_{IPW}=\frac1N\sum_{i=1}^N\frac{S_i\ell_i}{\pi_i},\qquad
 \E[\widehat R_{IPW}]=\frac1N\sum_{i=1}^N\ell_i
\end{equation}
利用的是$\E S_i=\pi_i$。查询越少的区域权重越大，故方差可能很高；确定性只选最高分时，未入选样本$\pi_i=0$，无法作此修正。自适应多轮查询的条件概率还依赖历史，最终入选概率不能随意替换为某轮的查询概率；同一批查询标签又参与训练时，固定评估器的无偏论证也不再直接适用。

一个可操作的补充是将部分预算留给概率探索：每轮候选数为$|U|$时，按$q_t(x)=(1-\epsilon)q_{a,t}(x)+\epsilon/|U|$查询，$0<\epsilon\leq1$，并记录当轮概率与历史状态。这保证当轮每个候选有非零机会，却不自动保证精确最终风险估计。最清晰的验收仍是保留不参与获取排序和训练的代表性监测样本，其成本应计入整体标注预算，而不是被当作免费信息。

标注噪声还会改变信息价值。若标注者$a$的混淆矩阵为$T^{(a)}_{ky}=P(\widetilde Y=k\mid Y=y,a)$，并假定给定真实标签后错误与$x,\theta$无关，则观测标签预测为$p(\widetilde y=k\mid x,\theta,a)=\sum_yT^{(a)}_{ky}p(y\mid x,\theta)$。查询信息增益应针对实际收到的$\widetilde Y$计算，而非假设完美的$Y$。矩阵本身未知、错误随图像质量变化或多人错误相关时，需要审核样本和更细模型；简单多数投票不能消除共同误判。

\paragraph{查询偏差与标注噪声的小例子。}
某记录以$\pi_i=1/4$概率被查询，则它被查询时在逆概率求和中贡献$4\ell_i$，未查询时贡献零，期望仍为$\ell_i$；除以总体数$N$后才是平均风险贡献。这不允许在$\pi_i=0$时补权重，也不意味着一次抽样结果一定接近真实平均。对于二分类标注者，若每类都有$0.1$概率被误标为另一类，则$T=\left(\begin{smallmatrix}0.9&0.1\\0.1&0.9\end{smallmatrix}\right)$，列对应真标签且每列和为一。模型真标签概率为$(0.8,0.2)^{\mathsf T}$时，收到标签的预测为$T(0.8,0.2)^{\mathsf T}=(0.74,0.26)^{\mathsf T}$，并非原概率。应把标注质量也纳入查询价值，不能只增加重复标注数量。

\section{算法流程与验收}
在巡检图像标注中，可以先按设备划分候选池和固定测试集，再从训练池取一批共同的初始样本。每轮用相同训练预算更新模型，给候选图像打获取分数，去掉连续拍摄的近重复图像后送专家标注，同时记录标签、耗时、费用和拒标原因。新标签加入训练集后再开始下一轮，不能把本轮查询结果提前用于本轮的分数计算。

比较随机、分层随机、不确定性和一种多样性方法，绘制真实测试召回随累计标注成本的曲线。重复不同种子集与查询随机种子，同时报告获取耗时、训练耗时、少数故障覆盖和失败查询比例。主动学习只有在同等总成本下提高真实风险表现时才有价值；节省标签却需要无法承受的反复重训，并非完整收益。

\paragraph{自测与答案。}
两个委员会成员都预测$(0.5,0.5)$，其预测熵和经验 BALD 分数各是多少？答案：预测熵为$\log2$，BALD 为$\log2-(\log2+\log2)/2=0$；高熵未必表示成员分歧。若本轮送标三条、两条有效返回、一条拒标，已标注集只增加两条，但三条都可能产生费用，不能把拒标记录当作已获得标签。

主动学习的获取函数、批量多样性和预算评价可用以下材料复习，分别对应算法直觉和工程实验。
\section*{辅助学习材料}
\begingroup
\emergencystretch=3em
\begin{itemize}
\item \textbf{Active Learning}；作者/平台：Small-Text；英文；官方教程。阅读 Components 和池式主动学习流程，理解初始化、查询策略、人工标注与停止准则如何组成迭代过程。\href{https://small-text.readthedocs.io/en/latest/active\_learning.html}{在线阅读}。
\item \textbf{Validation curves: plotting scores to evaluate models}；作者/平台：scikit-learn Developers；英文；官方示例文档。重点看训练集大小、验证曲线和重复评估，帮助把本章的累计标注预算、风险曲线与独立测试集评价落到实验设计。\href{https://scikit-learn.org/stable/modules/learning\_curve.html}{在线阅读}。
\item \textbf{通俗易懂讲AI--主动学习}；知乎；中文解说。沿查询、标注与更新循环理解预算分配；标题与主题据必应索引，原页受限。\href{https://zhuanlan.zhihu.com/p/676122996}{在线阅读}。
\item \textbf{用大白话讲主动学习（Active Learning）}；CSDN；中文博客。比较不确定性和样本多样性，评价必须使用独立测试集；标题与主题据必应索引，原页受限。\href{https://blog.csdn.net/qq\_51323826/article/details/144357831}{在线阅读}。
\item \textbf{主动学习(Active Learning，AL)的步骤、代码流程讲解及实战}；上传者：来包番茄沙司；bilibili；中文技术讲解。对照本章查询策略、不确定性度量与标注循环，检查每轮预算；评价仍使用独立测试集，不以查询样本上的准确率作为停止依据。2026年10月4日公开元数据播放14,386次，高于本次另一直接相关入门候选的1,501次；已核对公开流程简介、标题和计数，未逐帧观看。\href{https://www.bilibili.com/video/BV1y34y147xN/}{观看视频}。
\end{itemize}
\endgroup
\paragraph{本章小结。}
主动学习把标签视为需要分配的资源。不确定性寻找尚未确定的判断，BALD 进一步区分模型间分歧与单模型仍有的歧义，联合信息增益扣除批量冗余，预期风险方法直接考虑更新后的收益，而几何覆盖与梯度嵌入提供较便宜的批量代理。它们各自近似不同对象，因此获取得分不能跨方法直接解释为真实收益；大梯度不一定有益，高熵不一定可学，远距离也不一定对应新故障。

模型会判断错，专家也可能标错，同一批查询还可能重复，所以选样策略必须连同费用一起检验。上线效果应在独立且有代表性的数据上估计，不能只用主动挑出的难例测试。与持续学习结合时，可以分别替换选样方法和更新方法，查清收益来自哪里。如果多个工厂无法汇集原始记录，还要安排本地计算、更新通信与聚合。下一章讨论这种联邦协作；标签缺失、选择偏差和历史遗忘仍需各自处理。
